\begin{englishabstract}

Large Language Models (LLMs) have greatly improved natural language understanding,
making natural-language-driven Shell command generation increasingly practical.
However, in Chinese Shell scenarios, general-purpose LLMs still suffer from limited domain coverage,
command hallucinations, and high execution risks.

This thesis presents an integrated intelligent agent system for Chinese Shell interaction,
combining retrieval enhancement and safe execution.
The system adopts a client-server architecture with a locally quantized Qwen-7B model,
and introduces a hierarchical SQLite-based memory mechanism together with Bash/Zsh syntax adaptation
to improve multi-turn interaction stability.

For retrieval, a Retrieval-Augmented Generation module is built on ChromaDB and
BAAI/bge-small-zh-v1.5 embeddings.
The method combines vector similarity and keyword-hit signals for hybrid reranking,
and uses keyword routing to dispatch queries to five domain sub-collections
(commands, safety, tasks, patterns, and examples).
For execution, a three-tier safety framework is implemented, including permission grading,
a pre-execution pipeline (syntax check--risk assessment--effect simulation),
and structured execution logs.
In addition, a stepwise execution mechanism is designed for complex user requests,
supporting iterative generation, validation, confirmation, and execution for up to 12 steps.

Experimental results on 80 natural-language queries show that the proposed retrieval strategy
achieves Recall@1 of 0.90 and MRR of 0.9375, both outperforming a pure vector-retrieval baseline.

\englishkeyword{Large Language Model; Shell Agent; Retrieval-Augmented Generation; Hybrid Retrieval; Command Safety}
\end{englishabstract}
